\scanpage{36}
Considérons aussi la suite exacte transposée de (47)
\[
0\longrightarrow\mathcal O_S\xrightarrow{i'_1}\underset{\substack{\Vert\\ B}}{\check V}\xrightarrow{\pi'_1}\check L\longrightarrow0 \tag*{(51)}
\]
et notons tout de suite que, comme $V$ est
de rang 2, on a un isom canonique
\[
\check V\simeq V\otimes(\det V)^{-1} \tag*{(52)}
\]
déduit de l'accouplement \deleted{symétrique} \inserted{alterné}
\[
\begin{aligned} V\otimes V&\longrightarrow\det V\\ x\otimes y&\longmapsto x\wedge y\end{aligned} \tag*{(53)}
\]
d'ailleurs la suite exacte (47) définit un
isom. canonique
\[
\det V\xrightarrow[\sim]{\alpha}L \tag*{(54)}
\]
\deleted{en associant à} par la condition
\[
\alpha(i_1(\lambda)\wedge x)=\lambda\pi_1(x)\qquad\lambda\in\Gamma(-,L),\quad x\in\Gamma(-,V) \tag*{(55)}
\]
En conjuguant (52) et (54), on trouve un
iso can
\[
\check V\xleftarrow[\sim]{\alpha}V\otimes L^{-1} \tag*{(56)}
\]
\deleted{par la condition} associé à l'accouplement \uuline{alterné}
\[
\left\{\begin{aligned}
(V\otimes L^{-1})\otimes V&\xrightarrow{\varphi}\underline O_S\\[-2pt]
\underset{(V\otimes V)\otimes L^{-1}}{\unsure{\Vert|\wr|\simeq}}&
\end{aligned}\right. \tag*{(57)}
\]
défini par la condition que
\[
\varphi(i_1(\underset{\substack{\rotatebox{-90}{$\in$}\\\Gamma L}}{\lambda})\otimes\underset{\substack{\rotatebox{-90}{$\in$}\\\Gamma V}}{x}\otimes\underset{\substack{\rotatebox{-90}{$\in$}\\\GammaV}}{\mu})=\lambda\mu\,\pi_1(x) \tag*{(58)}
\]
% archive sheet number: 35
